\documentclass[italian]{article}
\usepackage{listings}
\usepackage{xcolor}
\usepackage{makeidx}
\makeindex

\title{
	Utility di compressione di immagini tramite DCT \\
	\normalsize Benchmark e comparatore della compressione DCT2
}

\date{A.A.: 2025/2026}

\author{
	\normalsize
	Nicolas Chines (matr. 899536) \\
	\normalsize
	Jacopo Borgato (matr. 866305)
}

\lstset{
	language=Python,
	basicstyle=\ttfamily\small,
	keywordstyle=\color{blue}\bfseries,
	commentstyle=\color{green}\itshape,
	stringstyle=\color{red},
	showstringspaces=false,
	breaklines=true,
	tabsize=4,
	numbers=left
}

\begin{document}
\maketitle


\centerline{\bfseries{Sommario}}
Il progetto consiste nella realizzazione di due software che implementano la compressione delle immagini basata su DCT.

Nella prima parte viene implementato l'algoritmo DCT2, confrontando successsivamente i tempi di esecuzione con quello implementato dalla libreria numpy.

Per la seconda parte il programma permette all'utente di scegliere da filesystem un'immagine bitmap in toni di grigio, applicare la compressione secondo i parametri specificati e visualizzare affiancate l'immagine originale e compressa.


\newpage

\section{Introduzione}
La funzione dct2 calcola la trasformata discreta del coseno 2D di una matrice 2D, tipicamente un'immagine. Questa trasformata converte i dati spaziali in componenti di frequenza. \`E una tecnica fondamentale per la compressione di immagini, alla base dello standard JPEG. 

L'operazione \`e separabile: equivale a applicare una DCT 1D prima sulle righe e poi sulle colonne (o viceversa). I coefficienti ottenuti permettono di ricostruire l'immagine originale tramite la trasformata inversa (idct2). Durante la compressione, i coefficienti ad alta frequenza, che risultano essere meno percepibili all'occhio umano, possono essere eliminati per ridurre la dimensione del file con una perdita di qualità limitata.

\section{Implementazione DCT / DCTN}
L'idea che sta alla base del calcolo della \texttt{DCT} \`e la trasformazione da base canonica in uno spazio vettoriale a quella dei coseni.\\

\subsection{Implementazione DCT1}
Come fattore di normalizzazione viene usato \texttt{1/sqrt(N)} per il primo coefficente, per i restanti invece \texttt{sqrt(2/N)}.\\
In fase di inizializzazione viene creato un vettore di output \texttt{y}. di lunghezza pari alla dimensione del vettore, pieno di zeri.

\begin{lstlisting}[language=Python]
N = len(v)
y = np.zeros(N, dtype=float)
\end{lstlisting}

Successivamente viene calcolato il coefficente \texttt{k} per ogni posizione nel seguente modo:


\begin{lstlisting}[language=Python]
if k == 0:
	Ck = 1.0 / math.sqrt(N)
else:
	Ck = math.sqrt(2.0 / N)
	
for i in range(N):
	ak = math.cos((math.pi * k * (2*i + 1)) / (2 * N)) * v[i]
	sum_cos += ak
	y[k] = Ck * sum_cos
\end{lstlisting}

Dove \texttt{Ck} \`e il fattore di normalizzazione, \texttt{ak} \`e la base della DCT, che rappresenta le diverse frequenze, la quale viene aggiunta alla somma dei coseni.\\
In questo modo non si ha pi\`u informazione locale come nella base canonica ma si ha informazione riguardante tutto il vettore.\\
Il risultato \`e quindi il vettore \texttt{y} di coefficienti che rappresenta l'input in termini di componenti di frequenza del coseno.

\subsection{Implementazione DCT2}

\begin{lstlisting}[language=Python]
# DCT1 for columns
for j in range(N):
	y_m[:,j] = dct(m_in[:,j])

# DCT1 for rows
for j in range(N):
	y_m[j,:] = np.transpose(dct(np.transpose(m_in[j,:])))
\end{lstlisting}

\printindex
\end{document}